%%%%%%%%%%  scan idx 214  |  folio 201  %%%%%%%%%%
\[
  \begin{aligned}
  P^0_{X/S}&\simeq\cO_X,\\
  P^1_{X/S}&\simeq\cO_X\oplus\Omega^1_{X/S},
  \end{aligned}
\]
où $\Omega^1_{X/S}$ est le faisceau des différentielles relatives de $X$ par rapport à $S$ où \emph{faisceau cotangent relatif}. On note $T_{X/S}$ le dual de $\Omega^1_{X/S}$, qu'on appelle \emph{faisceau tangent relatif}. Le faisceau $T_{X/S}$ (resp. $\Omega^1_{X/S}$) est le faisceau des sections d'un fibré vectoriel de rang fini sur $X$, qu'on appelle \emph{fibré tangent} (resp. \emph{cotangent}) \emph{relatif}, et qui est une $S$-variété mixte au-dessus de $X$.

L'homomorphisme $f^{-1}(\cO_S)$-linéaire (correspondant à la structure de $\cO_X$-Algèbre définie par $\pr_2$)
\[
  d^n_{X/S}:\cO_X\longrightarrow P^n_{X/S},
\]
qui à une section de $\cO_X$ associe sa partie principale à l'ordre $n$, est appelé \emph{opérateur différentiel universel d'ordre} $n$ sur $X$ relativement à $S$.

\textbf{1.2.2.} Plus généralement, si $E$ est un $\cO_X$-Module localement libre de type fini, on peut considérer le faisceau $P^n_{X/S}(E)=P^n_{X/S}\oplus_{\cO_X}E$ des parties principales à l'ordre $n$ de $E$ (le produit tensoriel étant défini au moyen de la deuxième structure de $\cO_X$-Module sur $P^n_{X/S}$), et l'on a un opérateur différentiel universel
\[
  d^n_{X/S}(E):E\longrightarrow P^n_{X/S}(E).
\]

Soient $E,F$ des $\cO_X$-Modules localement libres de type fini. Par définition, un opérateur différentiel de $E$ dans $F$ d'ordre $\le n$ est un homomorphisme $f^{-1}(\cO_X)$-linéaire de $E$ dans $F$ qui se factorise à travers $d^n_{X/S}(E)$ en un homomorphisme $\cO_X$-linéaire de $P^n_{X/S}(E)$ dans $F$ (la factorisation est alors unique). En d'autres termes, l'opérateur $d^n_{X/S}(E)$ établit une bijection
\begin{equation}\tag{1.2.2.1}
  \Homop_{\cO_X}\bigl(P^n_{X/S}(E),F\bigr)
  \xrightarrow{\sim}
  \Diff^n(E,F),
\end{equation}
où $\Diff^n(E,F)$ désigne l'ensemble des opérateurs différentiels de $E$
